% explainer-source-hash: f66797fe7a8c9d53245973007383b8fd1cc4dd9c9dfb65e80e4b5a7a95976308
% explainer-source-commit: a0ae5757644677d02ec29161881badb6a764a332
% explainer-generated: 2026-07-16
% Regenerate with /explain-all (see WIRING.md). Do not edit this stamp by hand:
% it is how the toolchain knows whether this document still matches the code.
\documentclass[11pt,a4paper]{article}
% preamble.tex - shared preamble for this project's explainer documents.
%
% This file is generated: it is copied in beside the documents on every build, so a
% clone of this repo can rebuild the PDFs, but local edits to it are overwritten.
%
% Usage from a document:
%   \documentclass[11pt,a4paper]{article}
%   % preamble.tex - shared preamble for this project's explainer documents.
%
% This file is generated: it is copied in beside the documents on every build, so a
% clone of this repo can rebuild the PDFs, but local edits to it are overwritten.
%
% Usage from a document:
%   \documentclass[11pt,a4paper]{article}
%   % preamble.tex - shared preamble for this project's explainer documents.
%
% This file is generated: it is copied in beside the documents on every build, so a
% clone of this repo can rebuild the PDFs, but local edits to it are overwritten.
%
% Usage from a document:
%   \documentclass[11pt,a4paper]{article}
%   \input{preamble}
%   \begin{document} ... \end{document}
%
% Build:  pdflatex -interaction=nonstopmode deep-dive.tex   (run it three times, so the
%         table of contents and the TikZ positioning settle)
%
% PACKAGE NOTES. Every package used here was checked present on the machine these
% documents were built on (Debian TeX Live 2023). If you rebuild elsewhere and a build
% fails, these are the likely gaps. Do not add a package without checking first:
%   kpsewhich <pkg>.sty
%
% Deliberately NOT used, because they were absent and the build dies without them:
%   algorithm2e, algorithm, algorithmic, algpseudocode  (texlive-science not installed)
%     -> the \begin{pseudocode} environment defined below stands in for them.
%   lmodern, charter, mathptmx, times, mathpazo, courier, newtx*, tgtermes
%     -> their .sty files RESOLVE but the font files themselves are absent, so the
%        build dies with "I can't find file `bchr8t'" or similar. Verified by
%        test-compiling each one. Only Computer Modern (the default) and helvet
%        actually render. kpsewhich reporting a font it cannot render is the trap:
%        do not "upgrade" the font stack without test-compiling it first.
%   amsmath, amssymb -> not loaded, so \text{} and \blacksquare are unavailable.
%
% A fuller font stack and algorithm2e would come from:
%   sudo apt install texlive-fonts-recommended texlive-science cm-super

\usepackage[T1]{fontenc}
\usepackage[utf8]{inputenc}
% Font: Computer Modern (default). See the note above before changing this.
\usepackage[a4paper,margin=2.4cm,headheight=14pt]{geometry}
% microtype WITHOUT font expansion. This is load-bearing, do not "restore" it:
% cm-super is not installed, so T1 Computer Modern resolves to BITMAP fonts, and
% pdfTeX's auto-expansion needs scalable ones. With expansion on, any document
% long enough to reach a second page dies with
%   "pdfTeX error (font expansion): auto expansion is only possible with scalable fonts"
% Protrusion (the bigger visual win) still works and stays on.
% `sudo apt install cm-super` would lift this; until then, leave expansion off.
\usepackage[protrusion=true,expansion=false]{microtype}
% Kill hyphen ligatures in TYPEWRITER text. This is load-bearing, not cosmetic. Under T1,
% \texttt{--verbose} silently renders as "-verbose", and \texttt{---} as "--", so a
% document can print a command-line flag as the WRONG flag and look perfectly fine doing
% it. Verified by rendering the page and reading the glyphs. lstlisting and \verb were
% always safe; this makes \code, \file and \texttt safe too, which is what prose uses.
% Ordinary prose is unaffected and still ligates, which is what you want there.
\DisableLigatures[-]{family=tt*}
\usepackage{graphicx}
\usepackage{xcolor}
\usepackage{booktabs}
\usepackage{enumitem}
\usepackage{fancyhdr}
\usepackage{titlesec}
\usepackage{tcolorbox}
\usepackage{listings}
\usepackage{float}
\usepackage{caption}
\usepackage{tikz}
\usepackage{forest}
\usepackage{pgfplots}
\pgfplotsset{compat=1.18}
\usetikzlibrary{arrows.meta,positioning,shapes.geometric,fit,backgrounds,calc,decorations.pathreplacing}
\tcbuselibrary{skins,breakable}
\usepackage[hidelinks,bookmarks=true,bookmarksnumbered=true]{hyperref}

% ===== palette ===========================================================
\definecolor{inkblue}{HTML}{1F3A5F}
\definecolor{accent}{HTML}{2E7DAF}
\definecolor{softgrey}{HTML}{F4F6F8}
\definecolor{linegrey}{HTML}{C8D0D8}
\definecolor{codebg}{HTML}{F7F7F4}
\definecolor{warnred}{HTML}{A33B3B}
\definecolor{okgreen}{HTML}{2E7D53}

% ===== headings ==========================================================
\titleformat{\section}{\Large\bfseries\color{inkblue}}{\thesection}{0.7em}{}
\titleformat{\subsection}{\large\bfseries\color{inkblue}}{\thesubsection}{0.6em}{}
\titleformat{\subsubsection}{\normalsize\bfseries\color{accent}}{\thesubsubsection}{0.5em}{}

\pagestyle{fancy}
\fancyhf{}
\fancyhead[L]{\small\color{inkblue}\nouppercase{\leftmark}}
\fancyhead[R]{\small\color{inkblue}\thepage}
\renewcommand{\headrulewidth}{0.4pt}
\renewcommand{\headrule}{\hbox to\headwidth{\color{linegrey}\leaders\hrule height \headrulewidth\hfill}}

% ===== code ==============================================================
\lstdefinestyle{house}{
  backgroundcolor=\color{codebg},
  basicstyle=\ttfamily\footnotesize,
  keywordstyle=\color{inkblue}\bfseries,
  commentstyle=\color{gray}\itshape,
  stringstyle=\color{okgreen},
  numbers=left, numberstyle=\tiny\color{gray}, numbersep=8pt,
  frame=single, rulecolor=\color{linegrey},
  breaklines=true, breakatwhitespace=false,
  showstringspaces=false, tabsize=2,
  captionpos=b, xleftmargin=14pt, framexleftmargin=14pt,
  columns=fullflexible, keepspaces=true,
  upquote=true,  % render ' and " as straight quotes, not curly
  % NON-ASCII INSIDE A LISTING IS A HARD BUILD FAILURE, not a cosmetic issue: listings
  % reads raw bytes, so inputenc dies on the first multi-byte character ("Invalid UTF-8
  % byte sequence") and the whole document fails to build. These mappings cover what
  % realistically turns up in source code: Latin-1 accents, Spanish punctuation, smart
  % quotes picked up from a word processor, the degree and section signs.
  % NOT COVERED, and it cannot be here: Urdu, Arabic and Balochi script. pdflatex has no
  % font for them. Describe or transliterate such text; never paste it into a listing.
  literate=%
    {á}{{\'a}}1 {é}{{\'e}}1 {í}{{\'i}}1 {ó}{{\'o}}1 {ú}{{\'u}}1
    {Á}{{\'A}}1 {É}{{\'E}}1 {Í}{{\'I}}1 {Ó}{{\'O}}1 {Ú}{{\'U}}1
    {à}{{\`a}}1 {è}{{\`e}}1 {ì}{{\`i}}1 {ò}{{\`o}}1 {ù}{{\`u}}1
    {ä}{{\"a}}1 {ë}{{\"e}}1 {ï}{{\"i}}1 {ö}{{\"o}}1 {ü}{{\"u}}1
    {Ä}{{\"A}}1 {Ö}{{\"O}}1 {Ü}{{\"U}}1 {ß}{{\ss}}1
    {ñ}{{\~n}}1 {Ñ}{{\~N}}1 {ç}{{\c c}}1 {Ç}{{\c C}}1
    {¿}{{\textquestiondown}}1 {¡}{{\textexclamdown}}1
    {°}{{\textdegree}}1 {§}{{\S}}1 {£}{{\pounds}}1 {€}{{EUR}}1
    {“}{{\textquotedblleft}}1 {”}{{\textquotedblright}}1
    {‘}{{\textquoteleft}}1 {’}{{\textquoteright}}1
    {–}{{-}}1 {—}{{-}}1 {…}{{...}}1 {→}{{$\rightarrow$}}1 {✓}{{x}}1
}
\lstset{style=house}

% ===== callout boxes =====================================================
% \begin{keyidea}{Title} ... \end{keyidea}   - a concept a beginner must grasp
\newtcolorbox{keyidea}[1]{
  enhanced, breakable, colback=softgrey, colframe=accent,
  boxrule=0.4pt, arc=2pt, left=8pt, right=8pt, top=6pt, bottom=6pt,
  title={\bfseries #1}, coltitle=white, colbacktitle=accent, fonttitle=\small
}

% \begin{jargon}{Term} definition \end{jargon}  - define a term at first use
\newtcolorbox{jargon}[1]{
  enhanced, breakable, colback=white, colframe=linegrey,
  boxrule=0.4pt, arc=2pt, left=8pt, right=8pt, top=6pt, bottom=6pt,
  title={\bfseries Jargon: #1}, coltitle=inkblue, colbacktitle=softgrey, fonttitle=\small
}

% \begin{honest}{Title} ... \end{honest}  - brutally honest finding (rule 18)
\newtcolorbox{honest}[1]{
  enhanced, breakable, colback=white, colframe=warnred,
  boxrule=0.6pt, arc=2pt, left=8pt, right=8pt, top=6pt, bottom=6pt,
  title={\bfseries #1}, coltitle=white, colbacktitle=warnred, fonttitle=\small
}

% \begin{pseudocode}{Caption} ... \end{pseudocode}
% Stands in for algorithm2e, which is NOT installed. Use \Step / \Indent inside.
\newtcolorbox{pseudocode}[1]{
  enhanced, breakable, colback=codebg, colframe=inkblue,
  boxrule=0.4pt, arc=2pt, left=10pt, right=8pt, top=6pt, bottom=6pt,
  fontupper=\ttfamily\small,
  title={\bfseries\rmfamily #1}, coltitle=white, colbacktitle=inkblue, fonttitle=\small
}
\newcommand{\Step}[1]{\par\noindent #1}
\newcommand{\Indent}[1]{\par\noindent\hspace*{2em}#1}

% ===== tikz house styles =================================================
\tikzset{
  box/.style   = {rectangle, rounded corners=2pt, draw=inkblue, fill=softgrey,
                  text=inkblue, align=center, inner sep=6pt, minimum height=9mm,
                  font=\small},
  extbox/.style= {rectangle, rounded corners=2pt, draw=linegrey, fill=white,
                  text=black!70, align=center, inner sep=6pt, minimum height=9mm,
                  font=\small, dashed},
  store/.style = {cylinder, shape border rotate=90, aspect=0.25, draw=accent,
                  fill=white, text=inkblue, align=center, inner sep=5pt,
                  font=\small},
  dec/.style   = {diamond, aspect=1.6, draw=accent, fill=white, text=inkblue,
                  align=center, inner sep=2pt, font=\scriptsize},
  flow/.style  = {-{Stealth[length=2.4mm]}, draw=inkblue, thick},
  dflow/.style = {-{Stealth[length=2.4mm]}, draw=accent, thick, dashed},
  % align=center is load-bearing: without it, a \\ inside an edge label throws
  % the misleading "perhaps a missing \item" error.
  lbl/.style   = {font=\scriptsize\itshape, text=black!65, fill=white,
                  inner sep=2pt, align=center}
}

% forest default for directory trees
\forestset{
  dirtree/.style={
    for tree={
      font=\ttfamily\small, grow'=0, child anchor=west, parent anchor=south,
      anchor=west, calign=first, inner sep=2pt,
      edge path={\noexpand\path [draw, \forestoption{edge}]
        (!u.south west) +(9pt,0) |- (.child anchor) \forestoption{edge label};},
      before typesetting nodes={if n=1{insert before={[,phantom]}}{}},
      fit=band, before computing xy={l=15pt}
    }
  }
}

% ===== misc ==============================================================
\captionsetup{font=small, labelfont={bf,color=inkblue}}
\setlist[itemize]{leftmargin=1.2em, itemsep=2pt, topsep=3pt}
\setlist[enumerate]{leftmargin=1.4em, itemsep=2pt, topsep=3pt}
\newcommand{\file}[1]{\texttt{\small #1}}
\newcommand{\code}[1]{\texttt{\small #1}}
\setcounter{tocdepth}{2}
\setcounter{secnumdepth}{3}

%   \begin{document} ... \end{document}
%
% Build:  pdflatex -interaction=nonstopmode deep-dive.tex   (run it three times, so the
%         table of contents and the TikZ positioning settle)
%
% PACKAGE NOTES. Every package used here was checked present on the machine these
% documents were built on (Debian TeX Live 2023). If you rebuild elsewhere and a build
% fails, these are the likely gaps. Do not add a package without checking first:
%   kpsewhich <pkg>.sty
%
% Deliberately NOT used, because they were absent and the build dies without them:
%   algorithm2e, algorithm, algorithmic, algpseudocode  (texlive-science not installed)
%     -> the \begin{pseudocode} environment defined below stands in for them.
%   lmodern, charter, mathptmx, times, mathpazo, courier, newtx*, tgtermes
%     -> their .sty files RESOLVE but the font files themselves are absent, so the
%        build dies with "I can't find file `bchr8t'" or similar. Verified by
%        test-compiling each one. Only Computer Modern (the default) and helvet
%        actually render. kpsewhich reporting a font it cannot render is the trap:
%        do not "upgrade" the font stack without test-compiling it first.
%   amsmath, amssymb -> not loaded, so \text{} and \blacksquare are unavailable.
%
% A fuller font stack and algorithm2e would come from:
%   sudo apt install texlive-fonts-recommended texlive-science cm-super

\usepackage[T1]{fontenc}
\usepackage[utf8]{inputenc}
% Font: Computer Modern (default). See the note above before changing this.
\usepackage[a4paper,margin=2.4cm,headheight=14pt]{geometry}
% microtype WITHOUT font expansion. This is load-bearing, do not "restore" it:
% cm-super is not installed, so T1 Computer Modern resolves to BITMAP fonts, and
% pdfTeX's auto-expansion needs scalable ones. With expansion on, any document
% long enough to reach a second page dies with
%   "pdfTeX error (font expansion): auto expansion is only possible with scalable fonts"
% Protrusion (the bigger visual win) still works and stays on.
% `sudo apt install cm-super` would lift this; until then, leave expansion off.
\usepackage[protrusion=true,expansion=false]{microtype}
% Kill hyphen ligatures in TYPEWRITER text. This is load-bearing, not cosmetic. Under T1,
% \texttt{--verbose} silently renders as "-verbose", and \texttt{---} as "--", so a
% document can print a command-line flag as the WRONG flag and look perfectly fine doing
% it. Verified by rendering the page and reading the glyphs. lstlisting and \verb were
% always safe; this makes \code, \file and \texttt safe too, which is what prose uses.
% Ordinary prose is unaffected and still ligates, which is what you want there.
\DisableLigatures[-]{family=tt*}
\usepackage{graphicx}
\usepackage{xcolor}
\usepackage{booktabs}
\usepackage{enumitem}
\usepackage{fancyhdr}
\usepackage{titlesec}
\usepackage{tcolorbox}
\usepackage{listings}
\usepackage{float}
\usepackage{caption}
\usepackage{tikz}
\usepackage{forest}
\usepackage{pgfplots}
\pgfplotsset{compat=1.18}
\usetikzlibrary{arrows.meta,positioning,shapes.geometric,fit,backgrounds,calc,decorations.pathreplacing}
\tcbuselibrary{skins,breakable}
\usepackage[hidelinks,bookmarks=true,bookmarksnumbered=true]{hyperref}

% ===== palette ===========================================================
\definecolor{inkblue}{HTML}{1F3A5F}
\definecolor{accent}{HTML}{2E7DAF}
\definecolor{softgrey}{HTML}{F4F6F8}
\definecolor{linegrey}{HTML}{C8D0D8}
\definecolor{codebg}{HTML}{F7F7F4}
\definecolor{warnred}{HTML}{A33B3B}
\definecolor{okgreen}{HTML}{2E7D53}

% ===== headings ==========================================================
\titleformat{\section}{\Large\bfseries\color{inkblue}}{\thesection}{0.7em}{}
\titleformat{\subsection}{\large\bfseries\color{inkblue}}{\thesubsection}{0.6em}{}
\titleformat{\subsubsection}{\normalsize\bfseries\color{accent}}{\thesubsubsection}{0.5em}{}

\pagestyle{fancy}
\fancyhf{}
\fancyhead[L]{\small\color{inkblue}\nouppercase{\leftmark}}
\fancyhead[R]{\small\color{inkblue}\thepage}
\renewcommand{\headrulewidth}{0.4pt}
\renewcommand{\headrule}{\hbox to\headwidth{\color{linegrey}\leaders\hrule height \headrulewidth\hfill}}

% ===== code ==============================================================
\lstdefinestyle{house}{
  backgroundcolor=\color{codebg},
  basicstyle=\ttfamily\footnotesize,
  keywordstyle=\color{inkblue}\bfseries,
  commentstyle=\color{gray}\itshape,
  stringstyle=\color{okgreen},
  numbers=left, numberstyle=\tiny\color{gray}, numbersep=8pt,
  frame=single, rulecolor=\color{linegrey},
  breaklines=true, breakatwhitespace=false,
  showstringspaces=false, tabsize=2,
  captionpos=b, xleftmargin=14pt, framexleftmargin=14pt,
  columns=fullflexible, keepspaces=true,
  upquote=true,  % render ' and " as straight quotes, not curly
  % NON-ASCII INSIDE A LISTING IS A HARD BUILD FAILURE, not a cosmetic issue: listings
  % reads raw bytes, so inputenc dies on the first multi-byte character ("Invalid UTF-8
  % byte sequence") and the whole document fails to build. These mappings cover what
  % realistically turns up in source code: Latin-1 accents, Spanish punctuation, smart
  % quotes picked up from a word processor, the degree and section signs.
  % NOT COVERED, and it cannot be here: Urdu, Arabic and Balochi script. pdflatex has no
  % font for them. Describe or transliterate such text; never paste it into a listing.
  literate=%
    {á}{{\'a}}1 {é}{{\'e}}1 {í}{{\'i}}1 {ó}{{\'o}}1 {ú}{{\'u}}1
    {Á}{{\'A}}1 {É}{{\'E}}1 {Í}{{\'I}}1 {Ó}{{\'O}}1 {Ú}{{\'U}}1
    {à}{{\`a}}1 {è}{{\`e}}1 {ì}{{\`i}}1 {ò}{{\`o}}1 {ù}{{\`u}}1
    {ä}{{\"a}}1 {ë}{{\"e}}1 {ï}{{\"i}}1 {ö}{{\"o}}1 {ü}{{\"u}}1
    {Ä}{{\"A}}1 {Ö}{{\"O}}1 {Ü}{{\"U}}1 {ß}{{\ss}}1
    {ñ}{{\~n}}1 {Ñ}{{\~N}}1 {ç}{{\c c}}1 {Ç}{{\c C}}1
    {¿}{{\textquestiondown}}1 {¡}{{\textexclamdown}}1
    {°}{{\textdegree}}1 {§}{{\S}}1 {£}{{\pounds}}1 {€}{{EUR}}1
    {“}{{\textquotedblleft}}1 {”}{{\textquotedblright}}1
    {‘}{{\textquoteleft}}1 {’}{{\textquoteright}}1
    {–}{{-}}1 {—}{{-}}1 {…}{{...}}1 {→}{{$\rightarrow$}}1 {✓}{{x}}1
}
\lstset{style=house}

% ===== callout boxes =====================================================
% \begin{keyidea}{Title} ... \end{keyidea}   - a concept a beginner must grasp
\newtcolorbox{keyidea}[1]{
  enhanced, breakable, colback=softgrey, colframe=accent,
  boxrule=0.4pt, arc=2pt, left=8pt, right=8pt, top=6pt, bottom=6pt,
  title={\bfseries #1}, coltitle=white, colbacktitle=accent, fonttitle=\small
}

% \begin{jargon}{Term} definition \end{jargon}  - define a term at first use
\newtcolorbox{jargon}[1]{
  enhanced, breakable, colback=white, colframe=linegrey,
  boxrule=0.4pt, arc=2pt, left=8pt, right=8pt, top=6pt, bottom=6pt,
  title={\bfseries Jargon: #1}, coltitle=inkblue, colbacktitle=softgrey, fonttitle=\small
}

% \begin{honest}{Title} ... \end{honest}  - brutally honest finding (rule 18)
\newtcolorbox{honest}[1]{
  enhanced, breakable, colback=white, colframe=warnred,
  boxrule=0.6pt, arc=2pt, left=8pt, right=8pt, top=6pt, bottom=6pt,
  title={\bfseries #1}, coltitle=white, colbacktitle=warnred, fonttitle=\small
}

% \begin{pseudocode}{Caption} ... \end{pseudocode}
% Stands in for algorithm2e, which is NOT installed. Use \Step / \Indent inside.
\newtcolorbox{pseudocode}[1]{
  enhanced, breakable, colback=codebg, colframe=inkblue,
  boxrule=0.4pt, arc=2pt, left=10pt, right=8pt, top=6pt, bottom=6pt,
  fontupper=\ttfamily\small,
  title={\bfseries\rmfamily #1}, coltitle=white, colbacktitle=inkblue, fonttitle=\small
}
\newcommand{\Step}[1]{\par\noindent #1}
\newcommand{\Indent}[1]{\par\noindent\hspace*{2em}#1}

% ===== tikz house styles =================================================
\tikzset{
  box/.style   = {rectangle, rounded corners=2pt, draw=inkblue, fill=softgrey,
                  text=inkblue, align=center, inner sep=6pt, minimum height=9mm,
                  font=\small},
  extbox/.style= {rectangle, rounded corners=2pt, draw=linegrey, fill=white,
                  text=black!70, align=center, inner sep=6pt, minimum height=9mm,
                  font=\small, dashed},
  store/.style = {cylinder, shape border rotate=90, aspect=0.25, draw=accent,
                  fill=white, text=inkblue, align=center, inner sep=5pt,
                  font=\small},
  dec/.style   = {diamond, aspect=1.6, draw=accent, fill=white, text=inkblue,
                  align=center, inner sep=2pt, font=\scriptsize},
  flow/.style  = {-{Stealth[length=2.4mm]}, draw=inkblue, thick},
  dflow/.style = {-{Stealth[length=2.4mm]}, draw=accent, thick, dashed},
  % align=center is load-bearing: without it, a \\ inside an edge label throws
  % the misleading "perhaps a missing \item" error.
  lbl/.style   = {font=\scriptsize\itshape, text=black!65, fill=white,
                  inner sep=2pt, align=center}
}

% forest default for directory trees
\forestset{
  dirtree/.style={
    for tree={
      font=\ttfamily\small, grow'=0, child anchor=west, parent anchor=south,
      anchor=west, calign=first, inner sep=2pt,
      edge path={\noexpand\path [draw, \forestoption{edge}]
        (!u.south west) +(9pt,0) |- (.child anchor) \forestoption{edge label};},
      before typesetting nodes={if n=1{insert before={[,phantom]}}{}},
      fit=band, before computing xy={l=15pt}
    }
  }
}

% ===== misc ==============================================================
\captionsetup{font=small, labelfont={bf,color=inkblue}}
\setlist[itemize]{leftmargin=1.2em, itemsep=2pt, topsep=3pt}
\setlist[enumerate]{leftmargin=1.4em, itemsep=2pt, topsep=3pt}
\newcommand{\file}[1]{\texttt{\small #1}}
\newcommand{\code}[1]{\texttt{\small #1}}
\setcounter{tocdepth}{2}
\setcounter{secnumdepth}{3}

%   \begin{document} ... \end{document}
%
% Build:  pdflatex -interaction=nonstopmode deep-dive.tex   (run it three times, so the
%         table of contents and the TikZ positioning settle)
%
% PACKAGE NOTES. Every package used here was checked present on the machine these
% documents were built on (Debian TeX Live 2023). If you rebuild elsewhere and a build
% fails, these are the likely gaps. Do not add a package without checking first:
%   kpsewhich <pkg>.sty
%
% Deliberately NOT used, because they were absent and the build dies without them:
%   algorithm2e, algorithm, algorithmic, algpseudocode  (texlive-science not installed)
%     -> the \begin{pseudocode} environment defined below stands in for them.
%   lmodern, charter, mathptmx, times, mathpazo, courier, newtx*, tgtermes
%     -> their .sty files RESOLVE but the font files themselves are absent, so the
%        build dies with "I can't find file `bchr8t'" or similar. Verified by
%        test-compiling each one. Only Computer Modern (the default) and helvet
%        actually render. kpsewhich reporting a font it cannot render is the trap:
%        do not "upgrade" the font stack without test-compiling it first.
%   amsmath, amssymb -> not loaded, so \text{} and \blacksquare are unavailable.
%
% A fuller font stack and algorithm2e would come from:
%   sudo apt install texlive-fonts-recommended texlive-science cm-super

\usepackage[T1]{fontenc}
\usepackage[utf8]{inputenc}
% Font: Computer Modern (default). See the note above before changing this.
\usepackage[a4paper,margin=2.4cm,headheight=14pt]{geometry}
% microtype WITHOUT font expansion. This is load-bearing, do not "restore" it:
% cm-super is not installed, so T1 Computer Modern resolves to BITMAP fonts, and
% pdfTeX's auto-expansion needs scalable ones. With expansion on, any document
% long enough to reach a second page dies with
%   "pdfTeX error (font expansion): auto expansion is only possible with scalable fonts"
% Protrusion (the bigger visual win) still works and stays on.
% `sudo apt install cm-super` would lift this; until then, leave expansion off.
\usepackage[protrusion=true,expansion=false]{microtype}
% Kill hyphen ligatures in TYPEWRITER text. This is load-bearing, not cosmetic. Under T1,
% \texttt{--verbose} silently renders as "-verbose", and \texttt{---} as "--", so a
% document can print a command-line flag as the WRONG flag and look perfectly fine doing
% it. Verified by rendering the page and reading the glyphs. lstlisting and \verb were
% always safe; this makes \code, \file and \texttt safe too, which is what prose uses.
% Ordinary prose is unaffected and still ligates, which is what you want there.
\DisableLigatures[-]{family=tt*}
\usepackage{graphicx}
\usepackage{xcolor}
\usepackage{booktabs}
\usepackage{enumitem}
\usepackage{fancyhdr}
\usepackage{titlesec}
\usepackage{tcolorbox}
\usepackage{listings}
\usepackage{float}
\usepackage{caption}
\usepackage{tikz}
\usepackage{forest}
\usepackage{pgfplots}
\pgfplotsset{compat=1.18}
\usetikzlibrary{arrows.meta,positioning,shapes.geometric,fit,backgrounds,calc,decorations.pathreplacing}
\tcbuselibrary{skins,breakable}
\usepackage[hidelinks,bookmarks=true,bookmarksnumbered=true]{hyperref}

% ===== palette ===========================================================
\definecolor{inkblue}{HTML}{1F3A5F}
\definecolor{accent}{HTML}{2E7DAF}
\definecolor{softgrey}{HTML}{F4F6F8}
\definecolor{linegrey}{HTML}{C8D0D8}
\definecolor{codebg}{HTML}{F7F7F4}
\definecolor{warnred}{HTML}{A33B3B}
\definecolor{okgreen}{HTML}{2E7D53}

% ===== headings ==========================================================
\titleformat{\section}{\Large\bfseries\color{inkblue}}{\thesection}{0.7em}{}
\titleformat{\subsection}{\large\bfseries\color{inkblue}}{\thesubsection}{0.6em}{}
\titleformat{\subsubsection}{\normalsize\bfseries\color{accent}}{\thesubsubsection}{0.5em}{}

\pagestyle{fancy}
\fancyhf{}
\fancyhead[L]{\small\color{inkblue}\nouppercase{\leftmark}}
\fancyhead[R]{\small\color{inkblue}\thepage}
\renewcommand{\headrulewidth}{0.4pt}
\renewcommand{\headrule}{\hbox to\headwidth{\color{linegrey}\leaders\hrule height \headrulewidth\hfill}}

% ===== code ==============================================================
\lstdefinestyle{house}{
  backgroundcolor=\color{codebg},
  basicstyle=\ttfamily\footnotesize,
  keywordstyle=\color{inkblue}\bfseries,
  commentstyle=\color{gray}\itshape,
  stringstyle=\color{okgreen},
  numbers=left, numberstyle=\tiny\color{gray}, numbersep=8pt,
  frame=single, rulecolor=\color{linegrey},
  breaklines=true, breakatwhitespace=false,
  showstringspaces=false, tabsize=2,
  captionpos=b, xleftmargin=14pt, framexleftmargin=14pt,
  columns=fullflexible, keepspaces=true,
  upquote=true,  % render ' and " as straight quotes, not curly
  % NON-ASCII INSIDE A LISTING IS A HARD BUILD FAILURE, not a cosmetic issue: listings
  % reads raw bytes, so inputenc dies on the first multi-byte character ("Invalid UTF-8
  % byte sequence") and the whole document fails to build. These mappings cover what
  % realistically turns up in source code: Latin-1 accents, Spanish punctuation, smart
  % quotes picked up from a word processor, the degree and section signs.
  % NOT COVERED, and it cannot be here: Urdu, Arabic and Balochi script. pdflatex has no
  % font for them. Describe or transliterate such text; never paste it into a listing.
  literate=%
    {á}{{\'a}}1 {é}{{\'e}}1 {í}{{\'i}}1 {ó}{{\'o}}1 {ú}{{\'u}}1
    {Á}{{\'A}}1 {É}{{\'E}}1 {Í}{{\'I}}1 {Ó}{{\'O}}1 {Ú}{{\'U}}1
    {à}{{\`a}}1 {è}{{\`e}}1 {ì}{{\`i}}1 {ò}{{\`o}}1 {ù}{{\`u}}1
    {ä}{{\"a}}1 {ë}{{\"e}}1 {ï}{{\"i}}1 {ö}{{\"o}}1 {ü}{{\"u}}1
    {Ä}{{\"A}}1 {Ö}{{\"O}}1 {Ü}{{\"U}}1 {ß}{{\ss}}1
    {ñ}{{\~n}}1 {Ñ}{{\~N}}1 {ç}{{\c c}}1 {Ç}{{\c C}}1
    {¿}{{\textquestiondown}}1 {¡}{{\textexclamdown}}1
    {°}{{\textdegree}}1 {§}{{\S}}1 {£}{{\pounds}}1 {€}{{EUR}}1
    {“}{{\textquotedblleft}}1 {”}{{\textquotedblright}}1
    {‘}{{\textquoteleft}}1 {’}{{\textquoteright}}1
    {–}{{-}}1 {—}{{-}}1 {…}{{...}}1 {→}{{$\rightarrow$}}1 {✓}{{x}}1
}
\lstset{style=house}

% ===== callout boxes =====================================================
% \begin{keyidea}{Title} ... \end{keyidea}   - a concept a beginner must grasp
\newtcolorbox{keyidea}[1]{
  enhanced, breakable, colback=softgrey, colframe=accent,
  boxrule=0.4pt, arc=2pt, left=8pt, right=8pt, top=6pt, bottom=6pt,
  title={\bfseries #1}, coltitle=white, colbacktitle=accent, fonttitle=\small
}

% \begin{jargon}{Term} definition \end{jargon}  - define a term at first use
\newtcolorbox{jargon}[1]{
  enhanced, breakable, colback=white, colframe=linegrey,
  boxrule=0.4pt, arc=2pt, left=8pt, right=8pt, top=6pt, bottom=6pt,
  title={\bfseries Jargon: #1}, coltitle=inkblue, colbacktitle=softgrey, fonttitle=\small
}

% \begin{honest}{Title} ... \end{honest}  - brutally honest finding (rule 18)
\newtcolorbox{honest}[1]{
  enhanced, breakable, colback=white, colframe=warnred,
  boxrule=0.6pt, arc=2pt, left=8pt, right=8pt, top=6pt, bottom=6pt,
  title={\bfseries #1}, coltitle=white, colbacktitle=warnred, fonttitle=\small
}

% \begin{pseudocode}{Caption} ... \end{pseudocode}
% Stands in for algorithm2e, which is NOT installed. Use \Step / \Indent inside.
\newtcolorbox{pseudocode}[1]{
  enhanced, breakable, colback=codebg, colframe=inkblue,
  boxrule=0.4pt, arc=2pt, left=10pt, right=8pt, top=6pt, bottom=6pt,
  fontupper=\ttfamily\small,
  title={\bfseries\rmfamily #1}, coltitle=white, colbacktitle=inkblue, fonttitle=\small
}
\newcommand{\Step}[1]{\par\noindent #1}
\newcommand{\Indent}[1]{\par\noindent\hspace*{2em}#1}

% ===== tikz house styles =================================================
\tikzset{
  box/.style   = {rectangle, rounded corners=2pt, draw=inkblue, fill=softgrey,
                  text=inkblue, align=center, inner sep=6pt, minimum height=9mm,
                  font=\small},
  extbox/.style= {rectangle, rounded corners=2pt, draw=linegrey, fill=white,
                  text=black!70, align=center, inner sep=6pt, minimum height=9mm,
                  font=\small, dashed},
  store/.style = {cylinder, shape border rotate=90, aspect=0.25, draw=accent,
                  fill=white, text=inkblue, align=center, inner sep=5pt,
                  font=\small},
  dec/.style   = {diamond, aspect=1.6, draw=accent, fill=white, text=inkblue,
                  align=center, inner sep=2pt, font=\scriptsize},
  flow/.style  = {-{Stealth[length=2.4mm]}, draw=inkblue, thick},
  dflow/.style = {-{Stealth[length=2.4mm]}, draw=accent, thick, dashed},
  % align=center is load-bearing: without it, a \\ inside an edge label throws
  % the misleading "perhaps a missing \item" error.
  lbl/.style   = {font=\scriptsize\itshape, text=black!65, fill=white,
                  inner sep=2pt, align=center}
}

% forest default for directory trees
\forestset{
  dirtree/.style={
    for tree={
      font=\ttfamily\small, grow'=0, child anchor=west, parent anchor=south,
      anchor=west, calign=first, inner sep=2pt,
      edge path={\noexpand\path [draw, \forestoption{edge}]
        (!u.south west) +(9pt,0) |- (.child anchor) \forestoption{edge label};},
      before typesetting nodes={if n=1{insert before={[,phantom]}}{}},
      fit=band, before computing xy={l=15pt}
    }
  }
}

% ===== misc ==============================================================
\captionsetup{font=small, labelfont={bf,color=inkblue}}
\setlist[itemize]{leftmargin=1.2em, itemsep=2pt, topsep=3pt}
\setlist[enumerate]{leftmargin=1.4em, itemsep=2pt, topsep=3pt}
\newcommand{\file}[1]{\texttt{\small #1}}
\newcommand{\code}[1]{\texttt{\small #1}}
\setcounter{tocdepth}{2}
\setcounter{secnumdepth}{3}


\begin{document}

\begin{titlepage}
\centering
\vspace*{2.5cm}
{\Huge\bfseries\color{inkblue} ZaraiLink: Demo\par}
\vspace{0.6cm}
{\Large\color{accent} Deep Dive\par}
\vspace{1.2cm}
{\large An exhaustive, beginner-safe walkthrough of the public demo\par}
\vspace{2.5cm}
\begin{minipage}{0.82\textwidth}
\begin{honest}{Read this before anything else}
\textbf{This document describes a demo, not a product.} Every company, contact,
phone number, email address, price and shipment figure in this repository is
\textbf{fabricated}. Nothing here is real trade data. No part of this code talks to a
network. It is a hand-built imitation of another application's screens.

\medskip
\textbf{ZaraiLink itself is a team project, not solo work.} It was built by three
people. Salman Adnan led it; Umar Kashif and Fahad Nadeem wrote real and substantial
parts of it. The real system is private and its source is not in this repository. This
demo exists so the product can be shown without publishing code that three people
would have to agree to publish.
\end{honest}
\end{minipage}
\vfill
{\small\color{inkblue}Working document for the owner's own understanding and interview defense.\par}
\end{titlepage}

\tableofcontents
\newpage

% =====================================================================
\section{What this repository is, and what it is not}

\subsection{The one-paragraph version}

This repository, \file{zarailink-demo}, is a \textbf{static web page}: three JavaScript
files, one stylesheet, one HTML file, a font and two favicons. Opening \file{index.html}
in a browser gives you a clickable imitation of six screens from a different, private
application called ZaraiLink. You can type a search, filter results, browse a company
directory, open a company profile, spend fake tokens to reveal fake phone numbers, and
star companies onto a watchlist. Every one of those actions is computed locally from a
hard-coded array of invented companies. There is no server, no database, no build step,
and no network request of any kind.

\begin{keyidea}{Why a demo exists at all}
The real ZaraiLink repository is private, and it cannot be made public by one person.
It has three authors, and publishing shared work is not a decision one author gets to
make alone. That leaves a problem: how do you show a hiring panel a product they cannot
be given the source to? The answer taken here is to rebuild the \emph{screens} against
\emph{invented} data. The visual design language and the interaction patterns are things
any browser loading the real app already downloads and renders. The private parts (the
natural-language pipeline, the ranking model, the customs data, the API routes) stay out
of this repository entirely.
\end{keyidea}

\subsection{The two framings that must never be dropped}

Both of the following are stated in the repository's own \file{README.md}, its
\file{index.html} banner and its footer, and both are load-bearing. If either is dropped
when this project is presented, the presentation becomes dishonest.

\begin{enumerate}
\item \textbf{The data is fabricated.} ``Meridian AgroTrade Ltd.'' does not exist.
  Neither does ``Elena Marsh'' or her phone number. The domains are all under
  \code{.example.com}, a reserved domain that can never resolve to a real site. The
  headline statistics (``50,000+ Trade Records'') are invented strings, not counts of
  anything.
\item \textbf{ZaraiLink is a team project led by Salman Adnan.} The demo shows features
  from a shared codebase. Some of those features were written by other people. This demo
  demonstrates that the features work as described; it does not claim their authorship.
\end{enumerate}

\begin{honest}{What this repository cannot prove, and should not be asked to}
This demo is evidence of \emph{UI reconstruction skill and product judgment}. It is
\textbf{not} evidence that the real ZaraiLink's backend works, because none of that
backend is here. A reader of this repository alone cannot verify that the natural-language
pipeline, the hybrid retrieval, or the learning-to-rank model exist or perform. Those
claims rest entirely on the owner's word and on the private repository.

The specific claim in \file{README.md} that Salman wrote ``64 of 97 commits'' is
\textbf{unverifiable from this repository}. The private repository's git history is not
here. That number is the owner's to stand behind, and an interviewer who asks for it
should be pointed at the private repository, not at this one.
\end{honest}

\subsection{The relationship between the two repositories}

\begin{figure}[H]
\centering
\resizebox{\textwidth}{!}{
\begin{tikzpicture}[node distance=8mm]
  \node[extbox, minimum width=52mm, minimum height=34mm, align=center] (real)
    {\textbf{ZaraiLink (the real system)}\\[3pt]
     \footnotesize PRIVATE repository\\
     \footnotesize Three authors, led by Salman\\[4pt]
     \footnotesize Django + DRF backend\\
     \footnotesize React frontend\\
     \footnotesize NLU pipeline\\
     \footnotesize Hybrid retrieval\\
     \footnotesize Learning-to-rank model\\
     \footnotesize Real Pakistani customs data};

  \node[box, minimum width=52mm, minimum height=34mm, right=42mm of real, align=center] (demo)
    {\textbf{zarailink-demo (this repo)}\\[3pt]
     \footnotesize PUBLIC\\
     \footnotesize Written by Salman alone\\[4pt]
     \footnotesize Plain HTML\\
     \footnotesize Plain CSS\\
     \footnotesize Vanilla JavaScript\\
     \footnotesize No backend, no build\\
     \footnotesize Fabricated data only};

  \draw[flow] ($(real.east)+(0,10mm)$) -- node[lbl, above, align=center]
    {screen structure,\\colour tokens,\\interaction patterns\\\textbf{copied}} ($(demo.west)+(0,10mm)$);

  \draw[dflow, draw=warnred] ($(real.east)+(0,-6mm)$) -- node[lbl, below, align=center, text=warnred]
    {backend logic, models,\\real data, API routes\\\textbf{deliberately NOT copied}} ($(demo.west)+(0,-6mm)$);

  \node[lbl, below=3mm of real, text width=52mm] {Source of truth. Not in this repo.};
  \node[lbl, below=3mm of demo, text width=52mm] {What you are reading about.};
\end{tikzpicture}}
\caption{The demo copies the visible surface and deliberately copies none of the private logic.}
\end{figure}

\begin{jargon}{Backend and frontend}
A \textbf{frontend} is the part of an application that runs in the user's browser and
draws what you see. A \textbf{backend} is the part that runs on a server somewhere else,
holds the database, and answers the frontend's questions over the network. The real
ZaraiLink has both. This demo has \emph{only} a frontend, and that frontend asks nobody
anything: it reads a JavaScript array sitting in the same folder.
\end{jargon}

% =====================================================================
\section{Layout of the repository}

\subsection{Every file, and why it exists}

\begin{figure}[H]
\centering
\begin{forest} dirtree
[zarailink-demo/
  [index.html {\rmfamily\itshape all six screens, as static markup}]
  [styles.css {\rmfamily\itshape colour tokens and component layout}]
  [demo-data.js {\rmfamily\itshape the fabricated companies and plans}]
  [icons.js {\rmfamily\itshape a hand-drawn SVG icon set}]
  [app.js {\rmfamily\itshape all behaviour: routing, filtering, unlocking}]
  [assets/
    [fonts/
      [inter-variable.woff2 {\rmfamily\itshape self-hosted font, 48 KB}]
    ]
  ]
  [favicon.svg {\rmfamily\itshape emerald tile with a Z}]
  [favicon.png {\rmfamily\itshape raster fallback}]
  [netlify.toml {\rmfamily\itshape two lines: publish this folder as-is}]
  [LICENSE {\rmfamily\itshape PolyForm Strict 1.0.0}]
  [README.md]
  [docs/
    [explainers/
      [deep-dive.tex {\rmfamily\itshape this document}]
      [brief.tex]
    ]
  ]
]
\end{forest}
\caption{The whole project. Five files carry all of the behaviour.}
\end{figure}

The entire application is roughly 1,100 lines across five files:

\begin{center}
\begin{tabular}{@{}llr@{}}
\toprule
\textbf{File} & \textbf{Role} & \textbf{Lines} \\
\midrule
\file{app.js}      & All behaviour and rendering & 399 \\
\file{styles.css}  & All presentation & 369 \\
\file{index.html}  & Static shell for six screens & 271 \\
\file{demo-data.js}& The fabricated dataset & 70 \\
\file{icons.js}    & SVG icon lookup table & 34 \\
\bottomrule
\end{tabular}
\end{center}

\subsection{How the browser loads it}

There is no build step. This deserves emphasis, because it is unusual for something that
imitates a React application.

\begin{jargon}{Build step}
Most modern web applications are written in a form the browser cannot run (React's JSX,
TypeScript, Sass), and a \textbf{build step} is a program that translates that source into
plain HTML, CSS and JavaScript before it is deployed. It is a whole extra machine to
install, configure, run and keep working. This demo skips it: what is in the repository is
exactly what the browser executes. You can double-click \file{index.html} and it works,
offline, forever, with no toolchain.
\end{jargon}

The three scripts are loaded in a fixed order at the bottom of \file{index.html}, and that
order matters:

\begin{lstlisting}
  <script src="demo-data.js"></script>
  <script src="icons.js"></script>
  <script src="app.js"></script>
\end{lstlisting}

\begin{figure}[H]
\centering
\resizebox{0.95\textwidth}{!}{
\begin{tikzpicture}[node distance=7mm]
  \node[box] (html) {\file{index.html}\\parsed top to bottom};
  \node[box, below=of html] (css) {\file{styles.css}\\\footnotesize plus \file{@font-face} loads Inter};
  \node[store, right=22mm of html] (data) {\file{demo-data.js}\\\footnotesize defines \code{COMPANIES},\\\footnotesize \code{tokenBalance}, \code{USER}, \code{STATS}};
  \node[box, below=of data] (icons) {\file{icons.js}\\\footnotesize defines \code{icon(name, size)}};
  \node[box, right=22mm of data] (app) {\file{app.js}\\\footnotesize defines everything else,\\\footnotesize then waits};
  \node[dec, below=12mm of app] (dom) {DOM\\ready?};
  \node[box, below=12mm of dom] (init) {Initial render:\\\footnotesize token pill, stats, plans,\\\footnotesize directory grid, listeners};

  \draw[flow] (html) -- (css);
  \draw[flow] (html) -- (data);
  \draw[flow] (data) -- (icons);
  \draw[flow] (icons.east) -| (app.south west);
  \draw[flow] (data) -- (app);
  \draw[flow] (app) -- (dom);
  \draw[flow] (dom) -- node[lbl, right] {\code{DOMContentLoaded}} (init);
\end{tikzpicture}}
\caption{Load order. \file{app.js} depends on globals that the two earlier scripts define.}
\end{figure}

\begin{honest}{The load order is an unguarded dependency}
\file{app.js} uses \code{COMPANIES}, \code{tokenBalance} and \code{icon()} as bare
globals. Nothing declares or checks that dependency. If someone reorders those three
\code{<script>} tags, or adds \code{defer} to only one of them, the page breaks with a
\code{ReferenceError} and no useful message. For a five-file demo that is an acceptable
trade (a module system would be more machinery than the problem deserves), but it is a
real coupling that lives only in the ordering of three lines of HTML, and it is worth
naming rather than pretending the files are independent.
\end{honest}

% =====================================================================
\section{The data: what was invented, and how it is shaped}

\subsection{Why the shape matters more than the values}

The values in \file{demo-data.js} are invented and worthless. The \emph{field names} are
the interesting part: they mirror the real application's components, so the demo exercises
the real product's actual information architecture rather than a simplified stand-in.

\begin{figure}[H]
\centering
\begin{forest} dirtree
[COMPANIES {\rmfamily\itshape array of 4}
  [id {\rmfamily\itshape "c1" .. "c4", used as the routing key}]
  [name / legal\_name {\rmfamily\itshape display name vs registered entity}]
  [country / district / province {\rmfamily\itshape drives the trade-direction pills}]
  [sector / type / established / employees {\rmfamily\itshape profile stat grid}]
  [verified / has\_key\_contacts {\rmfamily\itshape both needed for the Verified badge}]
  [website / last\_active {\rmfamily\itshape PRESENT BUT NEVER READ}]
  [products {\rmfamily\itshape array}
    [name / variety / value\_added]
  ]
  [avg\_price / volume / shipments / avg\_shipment {\rmfamily\itshape result-card stats}]
  [contacts {\rmfamily\itshape array, 1 to 2 per company}
    [id {\rmfamily\itshape "k1" .. "k5", the unlock key}]
    [name / designation]
    [is\_public {\rmfamily\itshape free, never locked}]
    [unlocked {\rmfamily\itshape MUTATED at runtime}]
    [phone / email / whatsapp {\rmfamily\itshape whatsapp may be null}]
  ]
]
\end{forest}
\caption{The company record. Two fields are decorative: nothing ever reads them.}
\end{figure}

The full fabricated dataset is small enough to state exactly:

\begin{center}
\begin{tabular}{@{}lllcc@{}}
\toprule
\textbf{Company} & \textbf{Country} & \textbf{Sector} & \textbf{Verified} & \textbf{Contacts} \\
\midrule
Meridian AgroTrade Ltd.     & Brazil   & Sweeteners  & yes & 2 \\
Cascade Commodities Group   & Brazil   & Sweeteners  & yes & 1 (public) \\
Sable Harvest Exports       & Brazil   & Sweeteners  & no  & 1 \\
Northgate Chemical Traders  & Pakistan & Fertilizers & yes & 1 \\
\bottomrule
\end{tabular}
\end{center}

Four companies, five contacts. One of those five (Priya Nair at Cascade) is marked
\code{is\_public: true}, so it starts unlocked and free. The other four cost one token
each. The user starts with \code{tokenBalance = 7}.

\begin{honest}{The dataset is smaller than the screens it has to fill}
Three of the four companies are Brazilian and sell dextrose; one is Pakistani and sells
urea. That is enough to make the country-based filtering visibly do \emph{something}, but
only just. As shown later, it means two of the four trade-direction pills return three
results and the other two return one result, and it means no screen ever shows enough rows
for the sidebar filters or the grid layout to be stress-tested. A panel clicking through
this will notice the world is thin.
\end{honest}

\subsection{Dead data}

Three declarations in \file{demo-data.js} are never read by anything. This was verified by
searching every file in the repository for each identifier.

\begin{center}
\begin{tabular}{@{}lll@{}}
\toprule
\textbf{Declaration} & \textbf{Looks like} & \textbf{Actually} \\
\midrule
\code{HS\_TREE} & a lookup of 3 HS codes & declared, referenced nowhere \\
\code{company.website} & a profile link & set on all 4, rendered never \\
\code{company.last\_active} & a recency signal & set on 3, rendered never \\
\bottomrule
\end{tabular}
\end{center}

\begin{jargon}{HS code}
The Harmonized System is an international numbering scheme for traded goods. Every
physical product that crosses a border has an HS code, and customs records are indexed by
it. \code{1702.3000} is dextrose. The real ZaraiLink lets a user search by this code, which
is why the demo's search box detects a leading digit and says so.
\end{jargon}

\code{HS\_TREE} is the most interesting of the three, because it is a real fragment of the
product's model (code, name, category) that got built and then never wired to a screen.
The Dashboard's ``4 Ways to Search'' card hard-codes the example \code{1702.3000} in the
HTML instead of reading it from \code{HS\_TREE}, and \code{runSearch()} hard-codes the
string \code{"HS 1702.3000: Dextrose"} rather than looking the code up. The data structure
that would make that dynamic exists and is ignored.

% =====================================================================
\section{The router: six screens in one file}

\subsection{How the screens work}

All six screens are present in \file{index.html} at once, as six \code{<section
class="view">} elements. CSS hides all of them, and shows only the one carrying an extra
class:

\begin{lstlisting}
.view { display: none; }
.view.is-active { display: block; }
\end{lstlisting}

Switching screens is therefore just moving one CSS class:

\begin{lstlisting}
function switchView(viewId) {
  document.querySelectorAll(".view").forEach(v =>
    v.classList.toggle("is-active", v.id === `view-${viewId}`));
  document.querySelectorAll(".zl-navbar__links > a").forEach(a =>
    a.classList.toggle("is-active", a.dataset.view === viewId));
  window.scrollTo({ top: 0, behavior: "smooth" });
}
\end{lstlisting}

The two-argument form of \code{classList.toggle} is the trick that makes this three lines
instead of ten: \code{toggle(name, condition)} adds the class when the condition is true
and removes it when false, so one pass over every view both shows the target and hides all
the others.

\begin{jargon}{The DOM, and \code{querySelectorAll}}
The \textbf{DOM} (Document Object Model) is the browser's live, in-memory tree of the page.
JavaScript changes what you see by changing that tree. \code{querySelectorAll(".view")}
hands back every element whose class is \code{view}, using the same selector syntax CSS
uses.
\end{jargon}

\begin{figure}[H]
\centering
\resizebox{\textwidth}{!}{
\begin{tikzpicture}[node distance=13mm]
  \node[box] (dash) {\textbf{dashboard}\\\footnotesize search home};
  \node[box, right=26mm of dash] (results) {\textbf{results}\\\footnotesize Data Dashboard};
  \node[box, below=of dash] (dir) {\textbf{directory}\\\footnotesize Trade Directory};
  \node[box, right=26mm of dir] (profile) {\textbf{profile}\\\footnotesize Company Profile};
  \node[box, below=of dir] (watch) {\textbf{watchlist}};
  \node[box, right=26mm of watch] (sub) {\textbf{subscription}};

  \draw[flow] (dash) -- node[lbl, above] {\code{runSearch()}} (results);
  \draw[flow] (dash) -- node[lbl, left, align=center] {nav /\\quick pill} (dir);
  \draw[flow] (dir) -- node[lbl, above] {\code{openProfile(id)}} (profile);
  \draw[flow] (results) -- node[lbl, right, align=center] {result card\\\code{data-open-profile}} (profile);
  \draw[flow] (profile.south west) -- node[lbl, below, align=center] {Back to\\Directory} (dir.east);
  \draw[flow] (dir) -- node[lbl, left] {nav} (watch);
  \draw[flow] (watch) -- node[lbl, above] {nav} (sub);
  \draw[flow, dashed] (sub.north) to[out=60, in=-60] node[lbl, right, align=center] {``Buy Tokens''\\from modal} (profile.south east);
  \draw[flow] (results.north west) to[out=140, in=40] node[lbl, above] {Dashboard chip} (dash.north east);
\end{tikzpicture}}
\caption{The navigation graph. Every edge is a class swap; no URL ever changes.}
\end{figure}

\begin{honest}{No URL routing, and what that costs}
\code{switchView} never touches the address bar. There is no \code{history.pushState}, no
hash fragment, no router. Three consequences, all real:
\begin{itemize}
\item \textbf{The browser Back button leaves the demo.} It does not go back a screen; it
  leaves the page entirely. The only way back is the in-page ``Back to Directory'' link.
\item \textbf{No screen is linkable.} You cannot send someone a URL that opens a company
  profile. For a portfolio demo whose entire job is to be shown to other people, that is a
  genuine limitation.
\item \textbf{Reload resets everything.} State lives in plain variables, so a refresh
  restores 7 tokens and an empty watchlist.
\end{itemize}
For a scripted walkthrough this is defensible and the simplicity is worth something. But
``I skipped the router deliberately'' is a much better answer than being surprised by the
question.
\end{honest}

\subsection{Application state}

The entire state of the application is six module-level variables:

\begin{lstlisting}
const watchedIds = new Set();   // starred company ids
let currentScope = "import";    // Import | Export        <- never read
let currentRole  = "Supplier";  // Supplier | Buyer
let currentProfileId = null;    // which company the profile shows
let ddUnlocked   = false;       // category paywall cleared?
let ddActivePill = null;        // which trade-direction pill
let pendingUnlock = null;       // { companyId, contactId } awaiting confirm
\end{lstlisting}

plus \code{tokenBalance}, which lives in \file{demo-data.js} and is mutated from
\file{app.js}. There is no state container, no framework, and no persistence.

\begin{honest}{\code{tokenBalance} is defined in the data file and mutated from the app file}
\code{demo-data.js} declares \code{let tokenBalance = 7} alongside the fabricated
companies, and \code{app.js} writes to it from three different places. Mixing ``the
fixtures'' and ``the mutable session state'' in the file named \file{demo-data.js} is a
small architectural smell: everything else in that file is inert data, and this one line is
not. It works, but the file name lies slightly about its contents.

Worse, the \code{unlocked} flag on each contact is mutated \emph{in place} inside
\code{COMPANIES}. The fabricated fixture and the live session state are the same object.
It is why a reload is the only way to re-lock a contact.
\end{honest}

\subsection{One listener for the whole application}

Rather than attaching a handler to every button, \file{app.js} attaches a single
\code{click} listener to \code{document.body} and works out what was clicked:

\begin{lstlisting}
document.body.addEventListener("click", (e) => {
  const viewLink = e.target.closest("[data-view]");
  if (viewLink) { /* navigate */ return; }
  const openProfileBtn = e.target.closest("[data-open-profile]");
  if (openProfileBtn) { /* open profile */ return; }
  const watchBtn = e.target.closest("[data-watch]");
  if (watchBtn) { /* toggle star */ return; }
  // ... and so on
});
\end{lstlisting}

\begin{jargon}{Event delegation}
When you click a button, the browser fires the event on the button, then on its parent,
then its parent's parent, all the way up to \code{<body>}. This is called
\textbf{bubbling}. \textbf{Event delegation} exploits it: listen once at the top, and use
\code{e.target.closest(selector)} to walk back \emph{down} the chain and find which
interesting element was actually hit. \code{closest()} starts at the clicked element and
climbs until it finds an ancestor matching the selector, which is what makes clicking an
SVG icon \emph{inside} a button still register as clicking the button.
\end{jargon}

This is the single best decision in the codebase, and it is worth being able to explain
why. Almost every element on these screens is created by writing an HTML string into
\code{innerHTML}. Doing that \textbf{destroys and recreates} the elements, which would
throw away any listener attached directly to them. With delegation, the listener lives on
\code{<body>}, which is never recreated, so re-rendering a grid a hundred times costs
nothing and breaks nothing. Without it, every render would have to re-attach every handler.

\begin{figure}[H]
\centering
\resizebox{\textwidth}{!}{
\begin{tikzpicture}
  % == left column: tests 1 to 4 ==
  \node[box] (click) {click anywhere on \code{<body>}};
  \node[dec, below=6mm of click]  (d1) {\code{data-view}?};
  \node[dec, below=10mm of d1]    (d2) {\code{data-open-}\\\code{profile}?};
  \node[dec, below=10mm of d2]    (d3) {\code{data-watch}?};
  \node[dec, below=10mm of d3]    (d4) {\code{data-unlock-}\\\code{contact}?};

  \node[box, right=30mm of d1] (a1) {\code{switchView()}\\\footnotesize + re-render target};
  \node[box, right=30mm of d2] (a2) {\code{openProfile(id)}};
  \node[box, right=30mm of d3] (a3) {toggle \code{watchedIds}\\\footnotesize + \code{renderDirectory()}};
  \node[box, right=30mm of d4] (a4) {\code{requestUnlock(id)}};

  \draw[flow] (click) -- (d1);
  \foreach \a/\b in {d1/d2, d2/d3, d3/d4}
    \draw[flow] (\a) -- node[lbl, left] {no} (\b);
  \foreach \d/\a in {d1/a1, d2/a2, d3/a3, d4/a4}
    \draw[flow] (\d) -- node[lbl, above] {yes} (\a);

  % == right column: tests 5 to 7, continuing the same chain ==
  \node[dec, right=54mm of a1] (d5) {\code{data-pill}?};
  \node[dec, below=10mm of d5] (d6) {\code{data-buy-plan}?};
  \node[dec, below=10mm of d6] (d7) {\code{\#intel-toggle}?};
  \node[box, below=8mm of d7]  (close) {close the dropdown\\\footnotesize (the fall-through)};

  \node[box, right=30mm of d5] (a5) {\code{selectPill(id)}};
  \node[box, right=30mm of d6] (a6) {grant tokens};
  \node[box, right=30mm of d7] (a7) {toggle dropdown};

  \foreach \a/\b in {d5/d6, d6/d7, d7/close}
    \draw[flow] (\a) -- node[lbl, left] {no} (\b);
  \foreach \d/\a in {d5/a5, d6/a6, d7/a7}
    \draw[flow] (\d) -- node[lbl, above] {yes} (\a);

  % == chain continues from d4's "no" into d5, entering from the west ==
  \draw[flow] (d4.south) -- ++(0,-9mm) node[lbl, below left] {no}
    -| ($(d5.west)+(-14mm,0)$) -- (d5.west);
\end{tikzpicture}}
\caption{The delegated dispatcher, read left column first then right. Each \code{yes} branch ends
in \code{return}, so it is an ordered chain, not a switch: the first matching test wins and
nothing below it is consulted.}
\end{figure}

% =====================================================================
\section{The Dashboard and the search-mode detector}

\subsection{What the screen does}

The Dashboard is the real application's post-login home. It carries a ``4 Ways to Search''
explainer card, a large search box, an Import/Export toggle, quick pills, three example
queries, a dark statistics strip and a features grid.

The one piece of real logic on it is the \textbf{mode detector}: as you type, a coloured
pill appears telling you how the query will be interpreted.

\begin{lstlisting}
function detectMode(query, scopeSelected) {
  if (/^\d/.test(query.trim())) return "hscode";
  if (scopeSelected && query.trim()) return "ai";
  return "product";
}
\end{lstlisting}

\begin{jargon}{Regular expression}
\code{/\^{}\textbackslash d/} is a \textbf{regular expression}: a tiny pattern language for
matching text. \code{\^{}} means ``at the very start'', and \code{\textbackslash d} means
``any digit''. So the whole pattern asks: does this string start with a digit? A query
starting with a number is assumed to be an HS code.
\end{jargon}

\subsection{A branch that can never be reached}

\begin{honest}{``Product Search'' mode is dead code, and the README claims otherwise}
\code{detectMode} takes a \code{scopeSelected} argument, and its third branch returns
\code{"product"} only when that argument is falsy. But \textbf{both} call sites pass the
literal \code{true}:

\begin{lstlisting}
const mode = detectMode(query, true);   // in renderModePill
const mode = detectMode(q, true);       // in runSearch
\end{lstlisting}

Walk the branches with \code{scopeSelected} pinned to \code{true}:
\begin{itemize}
\item query starts with a digit $\rightarrow$ \code{"hscode"}
\item query is non-empty and does not start with a digit $\rightarrow$ \code{"ai"}
\item query is empty $\rightarrow$ \code{"product"}, \emph{but} \code{renderModePill}
  returns early on an empty query and clears the pill, so nothing renders.
\end{itemize}
Therefore \code{"product"} is \textbf{unreachable in the running application}. The green
``Product Search'' pill, and its \code{.mode-pill.is-product} CSS rule, can never appear.
This was confirmed by executing the function directly rather than by reading it.

This matters because \file{README.md} states that typing a query ``flips a colored pill
between HS Code Mode, Product Search, and AI Query Mode''. Two of those three are real. The
third cannot happen. \textbf{The README overstates the feature.}

The cause is visible one line up: \code{currentScope} exists to hold whether Import or
Export is selected, and \code{detectMode} clearly wants to receive it. Instead the call
sites hard-code \code{true}. Passing \code{currentScope} would not fix it either, since
\code{currentScope} is initialised to \code{"import"} and is therefore \emph{always} truthy.
The real design presumably starts the scope unselected (\code{null}) and only then does the
distinction between product search and AI query mean anything.
\end{honest}

\begin{honest}{The Import/Export toggle is decorative}
\code{currentScope} is assigned in exactly one place and read in \textbf{none}:
\begin{lstlisting}
document.querySelectorAll(".scope-btn").forEach(b => b.addEventListener("click", () => {
  currentScope = b.dataset.scope;                              // written...
  document.querySelectorAll(".scope-btn").forEach(x =>
    x.classList.toggle("is-active", x === b));                 // ...only the class matters
}));
\end{lstlisting}
Clicking Import or Export moves a green highlight and changes nothing else. Nothing in
search, filtering or results consults it. The Dashboard card even tells the user ``AI
Query: Select Import/Export first'', implying the toggle gates a behaviour. It does not.
\end{honest}

\subsection{Running a search}

\begin{lstlisting}
function runSearch(query) {
  const q = (query || document.getElementById("query-input").value || "dextrose").trim();
  const mode = detectMode(q, true);
  ddUnlocked = false;
  ddActivePill = null;
  document.getElementById("dd-title").textContent =
    mode === "hscode" ? "HS 1702.3000: Dextrose" : q;
  document.getElementById("dd-sub").textContent =
    "Sweeteners \u00B7 fabricated example results for this demo";
  document.getElementById("dd-shipment-count").textContent =
    `${COMPANIES.length * 14} shipments (demo)`;
  renderFilters();
  renderPills();
  // ... prompt to pick a pill, then:
  switchView("results");
}
\end{lstlisting}

Note the fallback chain: an explicit argument, else whatever is in the box, else the
literal \code{"dextrose"}. Searching with an empty box still produces a results screen.

\begin{honest}{The search query does not filter anything}
\code{runSearch} computes \code{q}, and then uses it \textbf{only as a heading}. The result
set is never filtered by it. Search ``PVC Resin'' and you get the same four dextrose and
urea companies as searching ``dextrose'', under a heading that says ``PVC Resin''. Search
any number at all and the heading hard-codes ``HS 1702.3000: Dextrose'' regardless of which
number you typed.

The screen is honest at the level of a label (\code{dd-sub} reads ``fabricated example
results for this demo'') but the illusion is thin, and an interviewer who types two
different searches will see it immediately. This is defensible as a scripted demo, but only
if it is stated up front rather than discovered.
\end{honest}

\begin{honest}{The shipment count is an arbitrary formula}
\code{COMPANIES.length * 14} evaluates to \textbf{56}, and the chip reads ``56 shipments
(demo)''. The dataset's own \code{shipments} fields sum to \textbf{125}. So the demo
displays a headline count that contradicts its own data, computed by multiplying a company
count by a magic number. \code{COMPANIES.reduce((a, c) => a + c.shipments, 0)} would have
been both simpler to justify and internally consistent. The \code{(demo)} suffix carries a
lot of weight here.
\end{honest}

% =====================================================================
\section{The Data Dashboard: results, pills and the paywall}

\subsection{The sidebar filters}

\code{renderFilters()} builds the ``Refine by Product'' checkbox list by flattening every
product name across every company:

\begin{lstlisting}
const products = [...new Set(COMPANIES.flatMap(c => c.products.map(p => p.name)))];
\end{lstlisting}

\begin{jargon}{\code{flatMap} and \code{Set}}
\code{map} turns each item into something else. \code{flatMap} does the same but then
flattens one level of nesting, which is what you want when each company yields an
\emph{array} of products and you want one flat list. A \code{Set} is a collection that
silently discards duplicates, so wrapping the list in \code{new Set(...)} and spreading it
back out with \code{[...]} is the idiomatic JavaScript way to say ``unique values''.
\end{jargon}

That yields three filters: Dextrose Monohydrate, Dextrose Anhydrous, Urea 46\%.

\begin{honest}{The filter checkboxes are inert}
\code{renderFilters} renders real \code{<input type="checkbox">} elements with real,
correctly computed counts beside them, and \textbf{no listener is ever attached to them}.
Ticking one does nothing. They are convincing precisely because the counts are right, which
makes them the most misleading control on the page: a checkbox that visibly responds to the
mouse and changes nothing reads as broken rather than as scripted.
\end{honest}

\subsection{The four trade-direction pills}

\begin{lstlisting}
const DD_PILLS = [
  { id: "foreign-suppliers", label: "Foreign Suppliers", sub: "Import → Sellers",
    filter: c => c.country !== "Pakistan" },
  { id: "foreign-buyers",    label: "Foreign Buyers",    sub: "Export → Buyers",
    filter: c => c.country !== "Pakistan" },
  { id: "pk-buyers",         label: "Pakistani Buyers",  sub: "Import → Local Buyers",
    filter: c => c.country === "Pakistan" },
  { id: "pk-suppliers",      label: "Pakistani Suppliers", sub: "Export → Local Sellers",
    filter: c => c.country === "Pakistan" }
];
\end{lstlisting}

Storing the filter as a function on the data object is a neat touch: \code{selectPill} does
not need to know what any pill \emph{means}, it just calls
\code{COMPANIES.filter(pill.filter)}.

\begin{honest}{Four pills, two behaviours}
Read the predicates. ``Foreign Suppliers'' and ``Foreign Buyers'' have \textbf{identical}
filters. So do ``Pakistani Buyers'' and ``Pakistani Suppliers''. Verified by running them:

\begin{center}
\begin{tabular}{@{}llc@{}}
\toprule
\textbf{Pill} & \textbf{Predicate} & \textbf{Results} \\
\midrule
Foreign Suppliers   & \code{country !== "Pakistan"} & 3 \\
Foreign Buyers      & \code{country !== "Pakistan"} & 3 (the same 3) \\
Pakistani Buyers    & \code{country === "Pakistan"} & 1 \\
Pakistani Suppliers & \code{country === "Pakistan"} & 1 (the same 1) \\
\bottomrule
\end{tabular}
\end{center}

The supplier/buyer axis is \textbf{pure decoration}: only the country axis does anything.
Clicking ``Foreign Buyers'' after ``Foreign Suppliers'' shows the identical three
Brazilian companies, and the only visible difference is the result-count line, which reads
``3 Buyers found'' instead of ``3 Suppliers found'' because of this:
\begin{lstlisting}
${results.length} ${pill.label.includes("Buyer") ? "Buyers" : "Suppliers"} found
\end{lstlisting}
That is a string test on a display label to decide a noun, which also means it silently
depends on the label never being reworded.

The demo's \file{README.md} is careful here and says the pills ``filter the same fabricated
company set by country'', which is accurate. The screen is less careful than the README.
\end{honest}

\subsection{The category paywall, and a real bug}

When a pill is selected and \code{ddUnlocked} is false, a banner is rendered below the
results offering to unlock the whole category.

\begin{honest}{The paywall advertises 500 tokens and charges 5}
This is a genuine defect, not a design choice. The button label and the arithmetic
disagree by a factor of 100:

\begin{lstlisting}
<button ... id="unlock-category-btn">Unlock for 500 Tokens</button>
...
if (tokenBalance < 5) { showModal("modal-insufficient"); return; }
tokenBalance -= 5;
\end{lstlisting}

The user is told it costs 500. The code requires 5 and deducts 5. Since the user starts
with 7 tokens, the button \emph{works}, which is exactly why the bug survives: nothing
appears broken. You click a button marked 500, you have 7, and it succeeds.

Worse, the cost 5 is a \textbf{magic number written twice}, once in the guard and once in
the deduction, with the advertised price a third, unrelated literal in an HTML string. All
three should be one constant. Compare the per-contact flow, which consistently uses 1 in
the label, the guard and the deduction.

Which number is wrong is a judgment call. Against the plan grants (Starter gives 20 tokens
for \$29/mo), a 500-token category unlock is absurd: no plan except Enterprise could ever
afford one. So the \emph{label} is almost certainly the typo and 5 is the intended price.
Either way, one of the two must change.
\end{honest}

\subsection{The locked result cards}

\begin{lstlisting}
.dd-card.is-locked .dd-card__actions { filter: blur(2px); pointer-events: none; opacity: .6; }
\end{lstlisting}

Locked cards blur their action buttons and set \code{pointer-events: none}, which makes
the element ignore the mouse entirely.

\begin{honest}{The blur is cosmetic, and here that is fine, but know why}
A CSS blur hides nothing. The text is in the DOM, and anyone can read it in the developer
tools or by deleting one class. In the real ZaraiLink, a paywall implemented this way
would be a serious data leak: the paid-for contact details would be sitting in the page for
anyone who pressed F12.

In \emph{this} demo it is harmless, because every value behind the blur is fabricated and
there is nothing to protect. But the pattern is only safe here for a reason that does not
transfer, and the correct answer to ``is this how the real one works?'' is that a real
paywall must withhold the data at the server and never send it.
\end{honest}

% =====================================================================
\section{The Trade Directory and the Watchlist}

\subsection{Rendering the grid}

\begin{lstlisting}
function renderDirectory() {
  document.getElementById("td-title").textContent =
    currentRole === "Supplier" ? "Find Suppliers" : "Find Buyers";
  // ... lede and switch-button text likewise
  const query = (document.getElementById("td-search-input").value || "").toLowerCase();
  const results = COMPANIES.filter(c => c.name.toLowerCase().includes(query));
  document.getElementById("directory-grid").innerHTML =
    results.map(c => companyCardHtml(c)).join("")
    || `<p style="color:var(--text-2)">No suppliers found.</p>`;
}
\end{lstlisting}

The directory search box \emph{does} work, unlike the Dashboard's: it filters on company
name, case-insensitively, on every keystroke. Lowercasing both sides is the standard trick
for a case-insensitive contains-match. Note the \code{||} fallback: \code{[].join("")}
produces an empty string, which is falsy, so the empty-state message appears automatically.

\begin{honest}{Three defects in one small function}
\begin{enumerate}
\item \textbf{\code{currentRole} changes only nouns.} Switching to ``Find Buyers'' rewrites
  the heading, the lede and the button, and shows \textbf{the same four companies}. There is
  no buyer/supplier distinction in the data at all. The Trade Directory has exactly one
  mode wearing two hats.
\item \textbf{The empty state is hard-coded to ``No suppliers found.''} Search for
  nonsense while in Buyer mode and it still says ``suppliers''. The one string that was
  supposed to vary with \code{currentRole} is the one that does not.
\item \textbf{The search box is never cleared.} \code{renderDirectory} reads the input but
  nothing resets it. Type ``Northgate'', navigate to a profile, come back via ``Back to
  Directory'', and the directory silently still shows one company, because the stale query
  is still in the box driving the filter.
\end{enumerate}
\end{honest}

\subsection{The watchlist, and a real state bug}

The star button on each card carries \code{data-watch="c1"}, and the delegated handler
toggles membership of a \code{Set}:

\begin{lstlisting}
const watchBtn = e.target.closest("[data-watch]");
if (watchBtn) {
  const id = watchBtn.dataset.watch;
  if (watchedIds.has(id)) watchedIds.delete(id); else watchedIds.add(id);
  renderDirectory();      // <-- always the directory, never the watchlist
  return;
}
\end{lstlisting}

\begin{honest}{Un-starring from the Watchlist screen does nothing visible}
\code{companyCardHtml} is shared by both grids, so the Watchlist screen renders star
buttons too. But the handler \textbf{unconditionally} calls \code{renderDirectory()} and
never \code{renderWatchlist()}. So on the Watchlist screen:

\begin{enumerate}
\item You click the star to remove a company.
\item \code{watchedIds} correctly drops the id. The model is now right.
\item \code{renderDirectory()} re-renders a grid on a \emph{hidden} screen.
\item The Watchlist grid in front of you is never re-rendered, so the star stays gold and
  the company stays on the list.
\end{enumerate}

The user sees a dead button. Navigating away and back makes the change appear, which is the
signature of exactly this class of bug: the model updated, the view did not. The fix is one
line, re-rendering whichever grid is currently active, or simply calling both. This is the
most likely bug for an interviewer to trip over by accident, because starring things and
then visiting the watchlist to unstar one is the obvious way to explore the feature.
\end{honest}

\subsection{The star icon trick}

\begin{lstlisting}
${icon("star", 15).replace('fill="none"', watched ? 'fill="currentColor"' : 'fill="none"')}
\end{lstlisting}

\code{icon()} returns an SVG string whose \emph{first} \code{fill="none"} is on the
\code{<svg>} element itself, so replacing the first occurrence flips the star between
outlined and filled.

\begin{honest}{It works because of an accident of attribute order}
\code{String.replace} with a string pattern replaces \textbf{only the first match}. This
code depends on \code{fill="none"} appearing on the \code{<svg>} wrapper before it appears
anywhere in the path data. That is true today. Reorder the attributes in \code{icon()},
or add an icon whose path contains \code{fill="none"}, and this silently paints the wrong
element. A \code{fill} parameter on \code{icon()} would cost three lines and remove the
landmine. Note also that the false branch replaces \code{fill="none"} with
\code{fill="none"}, which is a no-op written out longhand.
\end{honest}

% =====================================================================
\section{The Company Profile: masking, seniority, and the unlock flow}

\subsection{Masking}

Locked contact details are shown mangled rather than hidden, so the user can see that
something real is there:

\begin{lstlisting}
function maskEmail(email) {
  if (!email) return "";
  const [user, domain] = email.split("@");
  return `${user[0]}\u2022\u2022\u2022\u2022\u2022@${domain}`;
}
function maskPhone(phone) {
  if (!phone) return "";
  return phone.slice(0, 3) + " \u2022\u2022\u2022\u2022-\u2022\u2022\u2022\u2022\u2022\u2022\u2022";
}
\end{lstlisting}

\begin{jargon}{Why this listing shows \code{\textbackslash u2022}}
The real source contains the literal bullet character U+2022, typed directly. This
document's typesetter has no glyph for it and cannot print it, so each bullet is shown
above as \code{\textbackslash u2022}, which is JavaScript's escape for exactly the same
character. The two forms are interchangeable in real code: the listing above is valid
JavaScript and behaves identically to what is in the file. Only the notation differs.
\end{jargon}

So \code{e.marsh@meridianagrotrade.example.com} renders as the first letter, five bullets,
then the full domain; and \code{+55 11 4820-1187} renders as \code{+55} followed by a
fixed bullet pattern. A CSS blur is layered on top via \code{.contact-row.is-masked}.

\begin{honest}{The mask leaks the domain, and the phone mask lies about length}
Two things worth noticing:
\begin{itemize}
\item \code{maskEmail} preserves the \textbf{entire domain}. Since the domain is derived
  from the company name, and you are looking at that company's profile, that is not much of
  a leak here. But it means the mask hides only the local part.
\item \code{maskPhone} emits a \textbf{fixed} pattern regardless of the real number's
  length. The real number is 15 characters; the mask is always the first 3 plus the same
  12-character shape. It therefore does not preserve length, which is arguably better for
  privacy and worse for realism.
\item \code{maskEmail} would produce the string \code{"undefined"} followed by bullets if
  handed an email with no local part, since \code{user[0]} on an empty string is
  \code{undefined} and template literals stringify it. The \code{if (!email)} guard catches
  \code{null} and \code{""} but not \code{"@domain.com"}. No such value exists in the
  fabricated data, so this cannot fire today.
\end{itemize}
The README claims these use ``the same masking pattern as the real \code{maskValue()}''.
That claim cannot be checked from this repository, since the real function is not here.
\textbf{Treat it as inferred.}
\end{honest}

\subsection{The seniority classifier}

Each contact gets a coloured badge derived from their job title by an ordered chain of
regular expressions:

\begin{lstlisting}
function getSeniorityBadge(designation) {
  const d = (designation || "").toLowerCase();
  if (/(ceo|chief|director|owner|partner|president|founder)/.test(d))
    return ["Decision Maker", "decision"];
  if (/procurement/.test(d)) return ["Procurement", "procurement"];
  if (/sales/.test(d))       return ["Sales", "sales"];
  if (/(manag|head)/.test(d))return ["Management", "management"];
  return ["Team Member", "team"];
}
\end{lstlisting}

Note \code{manag} rather than \code{manager}: the truncated stem deliberately catches
``Manager'', ``Managing'' and ``Management'' with one pattern. That is a nice touch.

Running it against all five fabricated contacts gives:

\begin{center}
\begin{tabular}{@{}lll@{}}
\toprule
\textbf{Designation} & \textbf{Badge} & \textbf{Matched on} \\
\midrule
Head of Export Sales & Sales          & \code{sales} (before \code{head} is tested) \\
Managing Director    & Decision Maker & \code{director} \\
Procurement Manager  & Procurement    & \code{procurement} (before \code{manag}) \\
Founder \& CEO       & Decision Maker & \code{founder} \\
Sales Director       & Decision Maker & \code{director} (before \code{sales}) \\
\bottomrule
\end{tabular}
\end{center}

\begin{honest}{First match wins, and the ordering produces two odd results}
The chain is ordered by intended priority, and for the most part that is correct: a
``Managing Director'' should read as a decision maker, not as management, and the ordering
delivers that.

But look at the two \code{sales} rows. ``Sales Director'' is badged \textbf{Decision
Maker}, never Sales, because \code{director} is tested first. Meanwhile ``Head of Export
Sales'' is badged \textbf{Sales}, not Management, because \code{sales} is tested before
\code{head}. So of the two contacts with ``Sales'' in the title, the \emph{more} senior one
loses the Sales badge and the less senior one keeps it. Both outcomes are arguably
defensible, but they are an emergent consequence of ordering rather than a decision anyone
made, and the classifier can only ever return one badge where a real title genuinely
carries two axes (function and seniority).

Of five fabricated contacts, three land on ``Decision Maker''. The ``Management'' and
``Team Member'' branches are never exercised by the shipped data, so two of the five CSS
badge colours are dead in practice.
\end{honest}

\subsection{The unlock flow}

This is the centrepiece of the demo, and the thing the README is proudest of: the real
application unlocks contacts \textbf{one at a time}, not per company, and the demo
reproduces that.

\begin{figure}[H]
\centering
\resizebox{0.98\textwidth}{!}{
\begin{tikzpicture}[node distance=9mm]
  \node[box] (btn) {User clicks\\``Unlock Contact Details''};
  \node[box, below=of btn] (req) {\code{requestUnlock(contactId)}\\\footnotesize sets \code{pendingUnlock = \{companyId, contactId\}}};
  \node[box, below=of req] (modal) {\code{modal-confirm} opens\\\footnotesize ``Unlock X's details for 1 token?''};
  \node[dec, below=11mm of modal] (choice) {confirm?};
  \node[box, left=22mm of choice] (cancel) {\code{pendingUnlock = null}\\\footnotesize nothing spent};
  \node[dec, below=11mm of choice] (funds) {\code{tokenBalance}\\$\geq 1$?};
  \node[box, right=26mm of funds] (poor) {\code{modal-insufficient}\\\footnotesize \code{pendingUnlock = null}};
  \node[box, below=11mm of funds] (spend) {\code{tokenBalance -= 1}\\\code{k.unlocked = true}};
  \node[box, below=of spend] (rerender) {\code{renderTokenBalance()}\\\code{renderContacts(c)}};
  \node[box, below=of rerender] (ok) {\code{modal-success}\\\footnotesize shows the revealed values};

  \draw[flow] (btn) -- (req);
  \draw[flow] (req) -- (modal);
  \draw[flow] (modal) -- (choice);
  \draw[flow] (choice) -- node[lbl, above] {no} (cancel);
  \draw[flow] (choice) -- node[lbl, right] {yes} (funds);
  \draw[flow] (funds) -- node[lbl, above] {no} (poor);
  \draw[flow] (funds) -- node[lbl, right] {yes} (spend);
  \draw[flow] (spend) -- (rerender);
  \draw[flow] (rerender) -- (ok);
\end{tikzpicture}}
\caption{The per-contact unlock. Only the confirmed contact flips; every other contact on the page stays locked.}
\end{figure}

The \code{pendingUnlock} variable is what makes the modal work. The modal's OK button is a
single, fixed element in the HTML; it cannot know which contact it is for. So
\code{requestUnlock} stashes the target, and \code{confirmUnlock} picks it back up:

\begin{lstlisting}
function confirmUnlock() {
  hideModal("modal-confirm");
  if (!pendingUnlock) return;
  if (tokenBalance < 1) { showModal("modal-insufficient"); pendingUnlock = null; return; }
  const c = COMPANIES.find(x => x.id === pendingUnlock.companyId);
  const k = c.contacts.find(x => x.id === pendingUnlock.contactId);
  tokenBalance -= 1;
  k.unlocked = true;
  renderTokenBalance();
  renderContacts(c);
  document.getElementById("success-text").textContent = `${k.name} \u00B7 ${k.phone} \u00B7 ${k.email}`;
  showModal("modal-success");
  pendingUnlock = null;
}
\end{lstlisting}

This function is careful in ways worth pointing out under questioning: it re-checks the
balance at confirm time rather than trusting the check made when the button was drawn; it
clears \code{pendingUnlock} on every exit path including the failure ones; and it mutates
exactly one contact's \code{unlocked} flag before re-rendering the whole contact list from
the model. The pattern (change the model, then re-render from the model) is the right one,
and it is the same pattern the watchlist bug above gets wrong.

\begin{honest}{Two small holes in an otherwise solid flow}
\begin{itemize}
\item \textbf{The confirm modal cannot be dismissed by clicking the backdrop or pressing
  Escape.} \code{.zl-modal-overlay} covers the screen and only \code{\#confirm-cancel}
  closes it. Every modal on the page has this property. It is a one-listener fix and a
  standard expectation.
\item \textbf{The success modal restates the data that is already on the page.} It writes
  name, phone and email into \code{success-text} immediately after
  \code{renderContacts} has already revealed exactly those values behind it. Harmless, but
  redundant.
\item \textbf{Unlocking is irreversible without a reload}, since \code{unlocked} is mutated
  on the fixture itself. Demoing the locked state twice in one session is impossible.
\end{itemize}
\end{honest}

\subsection{The token economy, end to end}

\begin{center}
\begin{tabular}{@{}lrl@{}}
\toprule
\textbf{Event} & \textbf{Tokens} & \textbf{Consistent?} \\
\midrule
Starting balance          & $+7$   & \\
Unlock one contact        & $-1$   & yes: label, guard and deduction all say 1 \\
Unlock a category         & $-5$   & \textbf{no}: the button says 500 \\
Buy Starter (\$29/mo)     & $+20$  & \\
Buy Growth (\$79/mo)      & $+75$  & \\
Buy Enterprise (\$249/mo) & $+300$ & \\
\bottomrule
\end{tabular}
\end{center}

With four lockable contacts in the whole dataset and a starting balance of 7, a user can
unlock \emph{everything} and still have 3 tokens left. The scarcity the screen is built to
dramatise never actually bites.

\begin{honest}{Buying a plan is an unconditional token faucet}
\begin{lstlisting}
const plan = SUBSCRIPTION_PLANS.find(p => p.id === buyPlanBtn.dataset.buyPlan);
tokenBalance += plan.tokens;
renderTokenBalance();
buyPlanBtn.innerHTML = `${icon("check-circle", 13)} Tokens added`;
setTimeout(() => { buyPlanBtn.textContent = "Get Plan"; }, 2000);
\end{lstlisting}
No payment, no state, no limit. Click ``Get Plan'' on Enterprise ten times and you have
3,000 tokens. There is no notion of a \emph{current} plan at all: the plan cards are a
button that adds a number. For a demo that is a reasonable shortcut, but the Subscription
screen is the one screen where nothing is really being modelled.

Also note the two-second \code{setTimeout} restores the label with \code{textContent} after
setting it with \code{innerHTML}. That is correct here, and mildly lucky: the icon SVG is
written in as markup and then wiped by a plain-text assignment, which happens to be exactly
what is wanted.
\end{honest}

\begin{honest}{The insufficient-funds modal has a dead button}
\code{\#insufficient-buy} is labelled ``Buy Tokens'' and carries
\code{data-view="subscription"}, so the delegated navigation handler fires and switches to
the Subscription screen. But it \emph{also} has its own listener, added later, that calls
\code{hideModal}. Both run. The result is correct by accident: the modal closes and the
view changes. The ordering works only because the delegated handler is on \code{<body>} and
therefore runs after the element's own listener during bubbling. It is fragile, and the
\code{data-view} attribute alone would have done the whole job if \code{hideModal} were
called from the navigation branch.
\end{honest}

% =====================================================================
\section{Presentation: icons, fonts and colour}

\subsection{The icon set}

\file{icons.js} is a single function over a lookup table of 20 SVG path fragments:

\begin{lstlisting}
function icon(name, size) {
  size = size || 16;
  const paths = { "search": '<circle cx="11" cy="11" r="8"/><path d="m21 21-4.3-4.3"/>', /* ... */ };
  const body = paths[name] || "";
  return `<svg width="${size}" height="${size}" viewBox="0 0 24 24" fill="none"
    stroke="currentColor" stroke-width="2" stroke-linecap="round"
    stroke-linejoin="round" aria-hidden="true">${body}</svg>`;
}
\end{lstlisting}

\begin{jargon}{SVG, stroke, and \code{currentColor}}
\textbf{SVG} draws pictures with maths rather than pixels, so an icon stays sharp at any
size. \code{stroke} is the colour of the line; \code{fill} is the colour inside it. These
icons are outlines, so they set \code{fill="none"} and stroke only. \code{currentColor} is
a keyword meaning ``whatever the CSS text colour is here'', which is why one icon
definition works on a dark navbar and a white card without being redefined.
\end{jargon}

Every icon shares Lucide's 24 by 24 viewBox, 2px stroke and round caps, which is the whole
point: the real application uses Lucide, and matching those three parameters is what makes
a hand-drawn icon read as belonging to the same family.

\begin{honest}{An unknown icon name fails silently, and a third of the set is unused}
\code{paths[name] || ""} means \code{icon("serch")} returns a valid, empty, invisible SVG.
No error, no warning, just a gap. A \code{console.warn} in the fallback would cost one line.

Of the 20 icons defined, several (\code{shield-check}, \code{filter}, \code{x},
\code{moon}, \code{chevron-right}, \code{arrow-right}, \code{zap}, \code{chevron-down},
\code{bar-chart}, \code{search}) are never called through \code{icon()}. Some are genuinely
unused; others exist twice, because \file{index.html} pastes the \emph{same} SVG markup
inline rather than calling the function. The moon, the zap and the chevron in the navbar are
all hand-copied into the HTML while also sitting in the lookup table. There is one icon set
maintained in two places.
\end{honest}

\subsection{The font, and a stale README}

\file{styles.css} self-hosts a variable font:

\begin{lstlisting}
@font-face {
  font-family: 'Inter';
  src: url('assets/fonts/inter-variable.woff2') format('woff2-variations'),
       url('assets/fonts/inter-variable.woff2') format('woff2');
  font-weight: 100 900;
  font-display: swap;
}
\end{lstlisting}

\begin{jargon}{Variable font, and \code{font-display: swap}}
An ordinary font ships one file per weight. A \textbf{variable font} ships one file that
can be rendered at any weight on a continuous axis, which is what \code{font-weight: 100
900} declares. \code{font-display: swap} tells the browser to draw text immediately in a
fallback font and swap Inter in when it arrives, so the page is never blank while waiting.
\end{jargon}

The real ZaraiLink loads Inter from Google Fonts. This demo self-hosts it instead, so the
page stays fully offline and makes no external request. Listing the same file twice with
two \code{format()} hints is a deliberate fallback for browsers that do not understand
\code{woff2-variations}.

\begin{honest}{\file{README.md} contradicts the code about the font}
The README's Architecture section states:

\begin{quote}\itshape
``The real app loads 'Inter' from Google Fonts; this demo stays fully self-contained and
offline, so it falls back to the closest system font instead of fetching that externally.''
\end{quote}

That is \textbf{no longer true}. The demo does not fall back to a system font: it ships
\file{assets/fonts/inter-variable.woff2} and uses the real Inter. The git history shows why.
The commit ``Self-host real Inter variable font and replace emoji with Lucide-style SVG
icons'' came \emph{after} the README paragraph was written, and the paragraph was never
updated. \file{styles.css}'s own header comment is correct and says the demo self-hosts.
So the repository now describes its own font strategy two different ways in two files, and
the README is the wrong one.

\textbf{The code wins. The README should be corrected.} This is exactly the kind of detail
a careful reviewer checks, and being caught claiming a system-font fallback while shipping
a 48 KB font file is a bad look for a project whose entire pitch is honesty about what is
real.
\end{honest}

The font file is 48 KB, which is small for a full variable Inter (the complete Latin range
typically runs several times that). It is therefore \emph{probably} a subset containing
only the glyphs this page needs. \textbf{This is inferred from the file size and has not
been confirmed by inspecting the font tables.}

\subsection{Colour tokens}

\file{styles.css} declares the palette as CSS custom properties, which its header comment
states were copied 1:1 from the real application's \file{index.css} \code{:root} block:

\begin{lstlisting}
:root {
  --primary: #10B981;      /* emerald: the brand colour */
  --primary-hover: #059669;
  --slate-900: #0F172A;    /* the navbar and stats strip */
  --text: #1A1A1A;
  --border: #E5E7EB;
  --r-card: 16px;  --r-btn: 8px;  --wrap: 1200px;
}
\end{lstlisting}

\begin{jargon}{CSS custom properties (variables)}
A name starting with two hyphens declares a reusable value. \code{var(--primary)} then
substitutes it anywhere. Change the one declaration and every use updates. These are
usually called \textbf{design tokens} when they encode a design system's decisions rather
than incidental values.
\end{jargon}

\begin{honest}{One token is referenced but never declared}
\code{.zl-user-menu\_\_name} reads \code{var(--slate-200, \#e2e8f0)}. There is no
\code{--slate-200} in the \code{:root} block, so the fallback always wins. It renders
correctly (that fallback is the right colour) but the variable is a fiction. The palette
declares \code{slate-900}, \code{800}, \code{700} and \code{400}, and \code{200} was
either forgotten or dropped.

Similarly the \code{.zl-theme-toggle} rule styles a moon button with hover states, and
\file{index.html} renders that button with an \code{aria-label} of ``Toggle dark mode'',
but \textbf{no JavaScript is attached to it and no dark theme exists in the stylesheet}.
Clicking it does nothing at all. Same for the two quick pills ``I want to buy'' and ``I
want to sell'', which have no \code{data-view} and no handler; the two beside them do
navigate, which makes the dead pair harder to spot, not easier. And the \code{.zl-user-menu}
is styled with \code{cursor: pointer} and a hover background, advertising a dropdown that
does not exist.

That is four visible controls (theme toggle, two quick pills, user menu) plus the sidebar
checkboxes and the scope toggle that respond to the mouse and do nothing. For a demo whose
purpose is to be clicked by strangers, dead controls are the highest-cost defect class in
the repository: each one is an invitation for an interviewer to find the seam.
\end{honest}

% =====================================================================
\section{Security, licensing and deployment}

\subsection{XSS, and why it does not matter here}

Nearly everything is rendered by building HTML strings and assigning them to
\code{innerHTML}. There is an \code{escapeHtml} helper, and it is used on most interpolated
values:

\begin{lstlisting}
function escapeHtml(str) {
  return String(str == null ? "" : str)
    .replace(/&/g, "&amp;").replace(/</g, "&lt;").replace(/>/g, "&gt;");
}
\end{lstlisting}

\begin{jargon}{XSS (cross-site scripting)}
If an application takes text from somewhere and drops it straight into the page as HTML,
then text containing \code{<script>} becomes a real, executing script. That is
\textbf{XSS}. The defence is escaping: turning \code{<} into \code{\&lt;} so the browser
draws it as a character instead of interpreting it as markup.
\end{jargon}

\begin{honest}{The escaping is incomplete, and safe only because the data is a constant}
\code{escapeHtml} handles \code{\&}, \code{<} and \code{>}, but \textbf{not} quotes.
Several interpolations land \emph{inside} HTML attributes, for example
\code{data-open-profile="\$\{c.id\}"}. A value containing a double quote would break out of
the attribute, and no amount of \code{<}-escaping stops it. Escaping \code{"} and
\code{'} is the standard completion of this function.

Several other values skip \code{escapeHtml} entirely (\code{c.shipments},
\code{p.tokens}, the seniority label, the contact-summary string), and \code{icon()} output
is injected raw by design.

\textbf{None of this is exploitable}, and it is important to be precise about why: every
string in this application is a hard-coded literal in \file{demo-data.js}. There is no
user-supplied data, no network input, no URL parameter, and no storage. The two text
inputs are read but never rendered back as HTML. The attack surface is empty.

So this is not a vulnerability. It is a habit that would become one the instant this code
met real data, which is precisely what the real ZaraiLink does. The right answer to ``is
this safe?'' is: yes here, for a reason that would not survive contact with a backend.
\end{honest}

A full scan of every JavaScript, HTML, CSS and configuration file in the repository for
API keys, tokens, passwords, bearer credentials, private keys and connection strings
returned \textbf{nothing}. There are no secrets in this repository, which is what you would
expect from a codebase with no network calls, but it was checked rather than assumed. Every
domain in the fabricated data sits under \code{.example.com}, which is reserved by RFC 2606
and can never resolve.

\subsection{Deployment}

\begin{lstlisting}
[build]
  publish = "."
\end{lstlisting}

The whole of \file{netlify.toml}. It declares no build command, because there is nothing to
build: Netlify serves the folder verbatim. This is the payoff for having skipped a
toolchain.

\file{index.html} also carries \code{<meta name="robots" content="noindex">}, which asks
search engines not to index the page. That is a deliberate and correct choice: a page full
of invented companies with invented contact details should not turn up in a search for
those names.

\subsection{Licence}

The repository is under \textbf{PolyForm Strict 1.0.0}. This is worth understanding
because it is unusual and it is not an open-source licence.

\begin{keyidea}{What PolyForm Strict actually permits}
It grants permission to \emph{read} the source and nothing else. No use, no modification,
no distribution, no derivative works. It is source-available, not open source.

For portfolio work this is a coherent choice: the point is to let people read the code and
judge it, not to donate it. It is especially coherent \emph{here}, where the underlying
product is a private team project. But it should not be described as open source, and the
copyright line reads ``Copyright 2026 Salman Adnan'' on a repository whose subject matter
belongs to three people. The demo's code is Salman's alone, so that is defensible; the
distinction between ``this demo's code'' and ``ZaraiLink'' is doing real work in that
line.
\end{keyidea}

% =====================================================================
\section{Consolidated findings}

Everything below was verified by reading the source and, where it involved logic,
by executing it.

\subsection{Real defects}

\begin{center}
\begin{tabular}{@{}p{5.4cm}p{4.6cm}p{4.6cm}@{}}
\toprule
\textbf{Defect} & \textbf{Where} & \textbf{Effect} \\
\midrule
Category unlock says 500 tokens, charges 5 & \file{app.js} \code{selectPill} & Label and code disagree 100-fold \\
Un-starring on the Watchlist re-renders the Directory & \file{app.js} \code{data-watch} branch & Button looks dead to the user \\
``Product Search'' mode unreachable & \code{detectMode(q, true)} at both call sites & A documented feature cannot occur \\
Empty state hard-codes ``No suppliers found.'' & \code{renderDirectory} & Wrong noun in Buyer mode \\
Directory search box never cleared & \code{renderDirectory} & Stale filter on return from a profile \\
Shipment chip shows 56, data sums to 125 & \code{COMPANIES.length * 14} & Self-contradicting headline \\
\code{escapeHtml} does not escape quotes & \code{escapeHtml} & Latent; not exploitable today \\
\bottomrule
\end{tabular}
\end{center}

\subsection{Dead code and dead controls}

\begin{center}
\begin{tabular}{@{}p{5.4cm}p{9.6cm}@{}}
\toprule
\textbf{Thing} & \textbf{Status} \\
\midrule
\code{HS\_TREE} & Declared in \file{demo-data.js}, referenced nowhere \\
\code{company.website}, \code{company.last\_active} & Populated, never rendered \\
\code{currentScope} & Written on every scope click, read nowhere \\
Import/Export toggle & Moves a highlight, changes no behaviour \\
Theme toggle (moon) & Styled and labelled, no listener, no dark theme \\
``I want to buy'' / ``I want to sell'' pills & No \code{data-view}, no handler \\
User menu & \code{cursor: pointer}, hover state, no dropdown \\
``Refine by Product'' checkboxes & Real counts, no listener \\
\code{--slate-200} & Used with a fallback, never declared \\
10 of 20 icons in \file{icons.js} & Never called; several duplicated inline in HTML \\
\code{.mode-pill.is-product} & Styled for a state that cannot occur \\
Management / Team Member badges & No fabricated contact reaches those branches \\
\bottomrule
\end{tabular}
\end{center}

\subsection{Where the README overstates the code}

\begin{enumerate}
\item \textbf{Font.} README says the demo falls back to a system font. It self-hosts Inter.
  The code is right; the README is stale.
\item \textbf{Search modes.} README says the pill flips between three modes. Only two are
  reachable.
\item \textbf{Trade-direction pills.} README is accurate (``filter by country''), but the
  screen presents four pills that behave as two.
\end{enumerate}

\subsection{What the code does well}

It is worth being equally precise about this, because the list is real:

\begin{itemize}
\item \textbf{Event delegation from \code{<body>}.} The correct architecture for an
  \code{innerHTML}-driven UI, and the reason re-rendering is free.
\item \textbf{Per-contact unlock state.} The README calls this out as a deliberate choice
  over a simpler per-company toggle, and the code delivers it. \code{pendingUnlock} plus a
  confirm-time re-check of the balance is the right shape for a two-step transaction.
\item \textbf{Change the model, then re-render from the model.} \code{confirmUnlock} does
  this correctly. The watchlist bug is a deviation from a pattern the codebase otherwise
  gets right.
\item \textbf{Filters stored as functions on the pill data.} \code{selectPill} stays
  generic.
\item \textbf{No build step.} A demo whose whole job is to still work in two years, opened
  from a file path by a stranger, is exactly the case where a toolchain is a liability.
\item \textbf{The honesty scaffolding is real and everywhere.} A yellow banner at the top
  of every screen, a footer disclaimer, \code{noindex}, \code{.example.com} domains, and
  ``(demo)'' and ``fabricated'' printed into the results screen itself. This is not an
  afterthought.
\end{itemize}

\subsection{The most likely interview questions}

\begin{enumerate}
\item \textit{``Show me the search working.''} Type two different queries. The results do
  not change. Lead with this rather than being caught by it.
\item \textit{``What does the Import/Export toggle do?''} Nothing. It is wired to a
  variable nothing reads.
\item \textit{``Why does this say 500 tokens and take 5?''} A bug. Fix it before the
  interview.
\item \textit{``Which parts of ZaraiLink did you write?''} The README says 64 of 97
  commits, and this repository cannot back that up. Have the private repository's history
  ready, and name Umar Kashif and Fahad Nadeem before being asked.
\item \textit{``Is this the real product?''} No. Every screen is a reconstruction and every
  value is invented.
\end{enumerate}

\begin{honest}{The highest-value hour of work on this repository}
Not a new feature. Three things, in order:
\begin{enumerate}
\item \textbf{Fix the 500/5 token label} and the watchlist re-render. Both are one-line
  fixes and both are things a clicking interviewer will find.
\item \textbf{Correct the README's font paragraph} and soften the three-modes claim. A
  project sold on honesty cannot afford a README that overstates its own code, however
  innocently it drifted.
\item \textbf{Either wire up or delete the dead controls.} The theme toggle, the two quick
  pills, the scope toggle and the sidebar checkboxes. Deleting them costs nothing and
  removes six chances for the demo to look broken rather than scripted. A demo with fewer
  controls that all work is strictly stronger than one with more controls where half are
  scenery.
\end{enumerate}
The README's own ``What I'd do differently'' proposes adding the Trade Intelligence
cluster (link prediction, node2vec, the GNN work). That is the right \emph{ambition} and
the wrong \emph{next step}: adding a sixth screen to a demo with six dead controls makes
the seam wider, not narrower.
\end{honest}

\end{document}
